% Môn: Vật lý
% Lớp: 12
% Chương: Chương II. Khí lí tưởng
% Bài: L12C2 Bài 10. Định luật Charles · Đồ thị đẳng nhiệt trên các hệ tọa độ
% ===== Câu 24 =====
% ID: L12C2B10-202-TN
% Mức: VD
\dangbt{Đồ thị đẳng nhiệt trên các hệ tọa độ}
\begin{ex}
% ID: L12C2B10-202-TN
% Mức: VD
Trong dạng “Bài toán tổng hợp về quá trình đẳng nhiệt”, Trên hệ trục p- $V$, đường đẳng nhiệt là
\choice
{\True một nhánh hypebol}
{đường thẳng qua gốc}
{đường thẳng song song trục p}
{parabol}
\loigiai{
Đáp án A. Từ $p=C/V$, đồ thị là một nhánh hypebol.
}
\end{ex}

% ===== Câu 46 =====
% ID: L12C2B10-203-DS
% Mức: TH
\begin{ex}
% ID: L12C2B10-203-DS
% Mức: TH
Trong dạng “Nhận biết quá trình đẳng nhiệt và định luật Boyle”, Mô tả đồ thị đẳng nhiệt trong hệ p-V.
\choiceTF

\loigiai{
Đường là nhánh hypebol chữ nhật; nhiệt độ càng cao, đường đẳng nhiệt càng xa gốc tọa độ.
}
\end{ex}

% ===== Câu 68 =====
% ID: L12C2B10-204-TN
% Mức: VD
\begin{ex}
% ID: L12C2B10-204-TN
% Mức: VD
Trong dạng “Khối khí chịu áp suất của cột chất lỏng”, Trên hệ trục p- $V$, đường đẳng nhiệt là
\choice
{\True một nhánh hypebol}
{đường thẳng qua gốc}
{đường thẳng song song trục p}
{parabol}
\loigiai{
Đáp án A. Từ $p=C/V$, đồ thị là một nhánh hypebol.
}
\end{ex}

% ===== Câu 70 =====
% ID: L12C2B10-205-DS
% Mức: TH
\begin{ex}
% ID: L12C2B10-205-DS
% Mức: TH
Trong dạng “Tính áp suất, thể tích trong quá trình đẳng nhiệt”, Mô tả đồ thị đẳng nhiệt trong hệ p-V.
\choiceTF

\loigiai{
Đường là nhánh hypebol chữ nhật; nhiệt độ càng cao, đường đẳng nhiệt càng xa gốc tọa độ.
}
\end{ex}

% ===== Câu 93 =====
% ID: L12C2B10-206-DS
% Mức: TH
\begin{ex}
% ID: L12C2B10-206-DS
% Mức: TH
Xét các phát biểu sau trong dạng “Đồ thị đẳng nhiệt trên các hệ tọa độ”.
\choiceTF
{\True Trong quá trình đẳng nhiệt của một lượng khí xác định — phát biểu đúng là: $pV$ không đổi.}
{Trong quá trình đẳng nhiệt của một lượng khí xác định — phát biểu $p/T$ không đổi là đúng.}
{\True Hệ thức đúng của định luật Boyle là — phát biểu đúng là: $p_1V_1=p_2V_2$.}
{Hệ thức đúng của định luật Boyle là — phát biểu $V_1/T_1=V_2/T_2$ là đúng.}
\loigiai{
\begin{itemchoice}
\item (Đúng) Trong quá trình đẳng nhiệt của một lượng khí xác định — phát biểu đúng là: $pV$ không đổi. (Vì): ịnh luật Boyle cho $pV=\text{hằng số}$ khi $T$ không đổi. b) Sai: nội dung không phù hợp. c) Đúng: Hai trạng thái đẳng nhiệt thỏa $p_1V_1=p_2V_2$. d) Sai: nội dung không phù hợp
\item (Sai) Trong quá trình đẳng nhiệt của một lượng khí xác định — phát biểu $p/T$ không đổi là đúng.
\item (Đúng) Hệ thức đúng của định luật Boyle là — phát biểu đúng là: $p_1V_1=p_2V_2$.
\item (Sai) Hệ thức đúng của định luật Boyle là — phát biểu $V_1/T_1=V_2/T_2$ là đúng.
\end{itemchoice}
}
\end{ex}

% ===== Câu 94 =====
% ID: L12C2B10-207-DS
% Mức: TH
\begin{ex}
% ID: L12C2B10-207-DS
% Mức: TH
Trong dạng “Đồ thị đẳng nhiệt trên các hệ tọa độ”, Mô tả đồ thị đẳng nhiệt trong hệ p-V.
\choiceTF

\loigiai{
Đường là nhánh hypebol chữ nhật; nhiệt độ càng cao, đường đẳng nhiệt càng xa gốc tọa độ.
}
\end{ex}

% ===== Câu 97 =====
% ID: L12C2B10-208-DS
% Mức: TH
\begin{ex}
% ID: L12C2B10-208-DS
% Mức: TH
Xét các phát biểu sau trong dạng “Đồ thị đẳng nhiệt trên các hệ tọa độ”.
\choiceTF
{\True Nén đẳng nhiệt để thể tích giảm một nửa thì áp suất — phát biểu đúng là: tăng gấp đôi.}
{Nén đẳng nhiệt để thể tích giảm một nửa thì áp suất — phát biểu giảm một nửa là đúng.}
{\True Đường đẳng nhiệt trong hệ p- $V$ là — phát biểu đúng là: nhánh hypebol.}
{Đường đẳng nhiệt trong hệ p- $V$ là — phát biểu đường tròn là đúng.}
\loigiai{
\begin{itemchoice}
\item (Đúng) Nén đẳng nhiệt để thể tích giảm một nửa thì áp suất — phát biểu đúng là: tăng gấp đôi. (Vì): p tỉ lệ nghịch $V$ nên $V$ giảm 2 lần thì p tăng 2 lần. b) Sai: nội dung không phù hợp. c) Đúng: Phương trình $p=C/V$ là một nhánh hypebol. d) Sai: nội dung không phù hợp
\item (Sai) Nén đẳng nhiệt để thể tích giảm một nửa thì áp suất — phát biểu giảm một nửa là đúng.
\item (Đúng) Đường đẳng nhiệt trong hệ p- $V$ là — phát biểu đúng là: nhánh hypebol.
\item (Sai) Đường đẳng nhiệt trong hệ p- $V$ là — phát biểu đường tròn là đúng.
\end{itemchoice}
}
\end{ex}

% ===== Câu 100 =====
% ID: L12C2B10-209-TN
% Mức: TH
\begin{ex}
% ID: L12C2B10-209-TN
% Mức: TH
Trong dạng “Nhận biết quá trình đẳng nhiệt và định luật Boyle”, Đường đẳng nhiệt trong hệ p- $V$ là
\choice
{\True nhánh hypebol}
{đường thẳng qua gốc}
{đường tròn}
{parabol}
\loigiai{
Đáp án A. Phương trình $p=C/V$ là một nhánh hypebol.
}
\end{ex}

% ===== Câu 101 =====
% ID: L12C2B10-210-DS
% Mức: VD
\begin{ex}
% ID: L12C2B10-210-DS
% Mức: VD
Xét các phát biểu sau trong dạng “Đồ thị đẳng nhiệt trên các hệ tọa độ”.
\choiceTF
{\True Trên đồ thị p- $V$, đường đẳng nhiệt có nhiệt độ cao hơn nằm — phát biểu đúng là: xa gốc tọa độ hơn.}
{Trên đồ thị p- $V$, đường đẳng nhiệt có nhiệt độ cao hơn nằm — phát biểu gần gốc hơn là đúng.}
{\True Khí có $p_1=1$ atm, $V_1=4$ lít, nén còn 2 lít. $p_2$ bằng — phát biểu đúng là: 2 atm.}
{Khí có $p_1=1$ atm, $V_1=4$ lít, nén còn 2 lít. $p_2$ bằng — phát biểu 4 atm là đúng.}
\loigiai{
\begin{itemchoice}
\item (Đúng) Trên đồ thị p- $V$, đường đẳng nhiệt có nhiệt độ cao hơn nằm — phát biểu đúng là: xa gốc tọa độ hơn. (Vì): Với cùng $V$, $T$ cao hơn cho p lớn hơn nên đường nằm xa gốc. b) Sai: nội dung không phù hợp. c) Đúng: $p_2=p_1V_1/V_2=2$ atm. d) Sai: nội dung không phù hợp
\item (Sai) Trên đồ thị p- $V$, đường đẳng nhiệt có nhiệt độ cao hơn nằm — phát biểu gần gốc hơn là đúng.
\item (Đúng) Khí có $p_1=1$ atm, $V_1=4$ lít, nén còn 2 lít. $p_2$ bằng — phát biểu đúng là: 2 atm.
\item (Sai) Khí có $p_1=1$ atm, $V_1=4$ lít, nén còn 2 lít. $p_2$ bằng — phát biểu 4 atm là đúng.
\end{itemchoice}
}
\end{ex}

% ===== Câu 104 =====
% ID: L12C2B10-211-TN
% Mức: TH
\begin{ex}
% ID: L12C2B10-211-TN
% Mức: TH
Trong dạng “Nhận biết quá trình đẳng nhiệt và định luật Boyle”, Trên đồ thị p- $V$, đường đẳng nhiệt có nhiệt độ cao hơn nằm
\choice
{\True xa gốc tọa độ hơn}
{gần gốc hơn}
{trùng trục $V$}
{trùng trục p}
\loigiai{
Đáp án A. Với cùng $V$, $T$ cao hơn cho p lớn hơn nên đường nằm xa gốc.
}
\end{ex}

% ===== Câu 105 =====
% ID: L12C2B10-212-DS
% Mức: VD
\begin{ex}
% ID: L12C2B10-212-DS
% Mức: VD
Xét các phát biểu sau trong dạng “Đồ thị đẳng nhiệt trên các hệ tọa độ”.
\choiceTF
{\True Trong quá trình đẳng nhiệt, khối lượng riêng của khí tỉ lệ — phát biểu đúng là: thuận với áp suất.}
{Trong quá trình đẳng nhiệt, khối lượng riêng của khí tỉ lệ — phát biểu nghịch với áp suất là đúng.}
{\True Bơm xe đạp nóng lên rõ rệt thì quá trình nén — phát biểu đúng là: không hoàn toàn đẳng nhiệt.}
{Bơm xe đạp nóng lên rõ rệt thì quá trình nén — phát biểu đẳng áp là đúng.}
\loigiai{
\begin{itemchoice}
\item (Đúng) Trong quá trình đẳng nhiệt, khối lượng riêng của khí tỉ lệ — phát biểu đúng là: thuận với áp suất. (Vì): Với khối lượng cố định, $\rho=m/V$ và pV không đổi nên $\rho\sim p$. b) Sai: nội dung không phù hợp. c) Đúng: Nhiệt độ thay đổi nên không thể áp dụng Boyle chính xác như quá trình đẳng nhiệt. d) Sai: nội dung không phù hợp
\item (Sai) Trong quá trình đẳng nhiệt, khối lượng riêng của khí tỉ lệ — phát biểu nghịch với áp suất là đúng.
\item (Đúng) Bơm xe đạp nóng lên rõ rệt thì quá trình nén — phát biểu đúng là: không hoàn toàn đẳng nhiệt.
\item (Sai) Bơm xe đạp nóng lên rõ rệt thì quá trình nén — phát biểu đẳng áp là đúng.
\end{itemchoice}
}
\end{ex}

% ===== Câu 109 =====
% ID: L12C2B10-213-DS
% Mức: VDC
\begin{ex}
% ID: L12C2B10-213-DS
% Mức: VDC
Xét các phát biểu sau trong dạng “Đồ thị đẳng nhiệt trên các hệ tọa độ”.
\choiceTF
{\True Để quá trình gần đẳng nhiệt nên nén khí — phát biểu đúng là: chậm để trao đổi nhiệt.}
{Để quá trình gần đẳng nhiệt nên nén khí — phát biểu thật nhanh là đúng.}
{\True Áp suất tuyệt đối của cột khí bị giam bằng — phát biểu đúng là: áp suất khí quyển cộng hoặc trừ áp suất cột chất lỏng tùy bố trí.}
{Áp suất tuyệt đối của cột khí bị giam bằng — phát biểu chỉ bằng áp suất cột chất lỏng là đúng.}
\loigiai{
\begin{itemchoice}
\item (Đúng) Để quá trình gần đẳng nhiệt nên nén khí — phát biểu đúng là: chậm để trao đổi nhiệt. (Vì): Nén chậm giúp khí trao đổi nhiệt và giữ nhiệt độ gần không đổi. b) Sai: nội dung không phù hợp. c) Đúng: Cân bằng áp suất phải xét cả khí quyển và chênh áp thủy tĩnh. d) Sai: nội dung không phù hợp
\item (Sai) Để quá trình gần đẳng nhiệt nên nén khí — phát biểu thật nhanh là đúng.
\item (Đúng) Áp suất tuyệt đối của cột khí bị giam bằng — phát biểu đúng là: áp suất khí quyển cộng hoặc trừ áp suất cột chất lỏng tùy bố trí.
\item (Sai) Áp suất tuyệt đối của cột khí bị giam bằng — phát biểu chỉ bằng áp suất cột chất lỏng là đúng.
\end{itemchoice}
}
\end{ex}

% ===== Câu 137 =====
% ID: L12C2B10-214-TN
% Mức: TH
\begin{ex}
% ID: L12C2B10-214-TN
% Mức: TH
Trong dạng “Tính áp suất, thể tích trong quá trình đẳng nhiệt”, Đường đẳng nhiệt trong hệ p- $V$ là
\choice
{\True nhánh hypebol}
{đường thẳng qua gốc}
{đường tròn}
{parabol}
\loigiai{
Đáp án A. Phương trình $p=C/V$ là một nhánh hypebol.
}
\end{ex}

% ===== Câu 139 =====
% ID: L12C2B10-215-TN
% Mức: TH
\begin{ex}
% ID: L12C2B10-215-TN
% Mức: TH
Trong dạng “Tính áp suất, thể tích trong quá trình đẳng nhiệt”, Trên đồ thị p- $V$, đường đẳng nhiệt có nhiệt độ cao hơn nằm
\choice
{\True xa gốc tọa độ hơn}
{gần gốc hơn}
{trùng trục $V$}
{trùng trục p}
\loigiai{
Đáp án A. Với cùng $V$, $T$ cao hơn cho p lớn hơn nên đường nằm xa gốc.
}
\end{ex}

% ===== Câu 148 =====
% ID: L12C2B10-216-TN
% Mức: TH
\begin{ex}
% ID: L12C2B10-216-TN
% Mức: TH
Trong dạng “Đồ thị đẳng nhiệt trên các hệ tọa độ”, Đường đẳng nhiệt trong hệ p- $V$ là
\choice
{\True nhánh hypebol}
{đường thẳng qua gốc}
{đường tròn}
{parabol}
\loigiai{
Đáp án A. Phương trình $p=C/V$ là một nhánh hypebol.
}
\end{ex}

% ===== Câu 149 =====
% ID: L12C2B10-217-TN
% Mức: TH
\begin{ex}
% ID: L12C2B10-217-TN
% Mức: TH
Trong dạng “Đồ thị đẳng nhiệt trên các hệ tọa độ”, Trên đồ thị p- $V$, đường đẳng nhiệt có nhiệt độ cao hơn nằm
\choice
{\True xa gốc tọa độ hơn}
{gần gốc hơn}
{trùng trục $V$}
{trùng trục p}
\loigiai{
Đáp án A. Với cùng $V$, $T$ cao hơn cho p lớn hơn nên đường nằm xa gốc.
}
\end{ex}
